\section{Стандартная библиотека: основные концепции и линейные контейнеры}

\subsection{Что такое UB (неопределённое поведение) в C++ и какую опасность оно представляет? Приведите примеры.}
\begin{definition}[UB]
\textit{Undefined Behavior (UB)} — это ситуация, при которой стандарт C++ не задаёт никакого поведения программы. После UB компилятор вправе делать что угодно: программа может работать «как будто нормально», может аварийно завершиться, а может начать вести себя непредсказуемо.
\end{definition}

Опасность UB в том, что ошибка может проявляться не всегда и не в том месте, где она допущена. Из-за этого UB особенно трудно отлаживать.

\begin{example}
Примеры UB:
\begin{verbatim}
int* p = nullptr;
*p = 10;                 // разыменование nullptr

int a[3] = {1, 2, 3};
std::cout << a[10];      // выход за границы массива

int x;
std::cout << x;          // использование неинициализированной переменной

int y = 1;
int z = y / 0;           // деление на ноль для целых типов
\end{verbatim}
\end{example}

\begin{remark}
UB — это не ошибка компилятора и не исключение. Это нарушение правил языка, после которого поведение программы уже не гарантируется.
\end{remark}

\subsection{Как и с помощью инструкций \texttt{typedef} и \texttt{using} определяются псевдонимы типов? Для чего это применяется?}
\begin{definition}[Псевдоним типа]
Псевдоним типа — это новое имя для уже существующего типа.
\end{definition}

Классический синтаксис:
\begin{verbatim}
typedef unsigned long ulong;
\end{verbatim}

Современный синтаксис:
\begin{verbatim}
using ulong = unsigned long;
\end{verbatim}

\begin{itemize}
    \item \texttt{typedef} — старый способ;
    \item \texttt{using} — более читаемый и удобный, особенно для шаблонов.
\end{itemize}

Применение:
\begin{itemize}
    \item упрощение длинных типов;
    \item повышение читаемости;
    \item скрытие деталей реализации;
    \item облегчение сопровождения кода.
\end{itemize}

\subsection{Какие основные компоненты входят в состав стандартной библиотеки?}
\begin{definition}[Стандартная библиотека]
Стандартная библиотека C++ — это набор готовых компонентов: контейнеров, алгоритмов, итераторов, строк, потоков, умных указателей, математических и служебных средств.
\end{definition}

Основные блоки:
\begin{itemize}
    \item контейнеры (\texttt{std::vector}, \texttt{std::array}, \texttt{std::list}, \texttt{std::map} и др.);
    \item алгоритмы (\texttt{std::sort}, \texttt{std::find}, \texttt{std::count});
    \item итераторы;
    \item строки и строковые потоки;
    \item потоки ввода-вывода;
    \item служебные заголовки и утилиты.
\end{itemize}

\subsection{Что такое контейнеры, алгоритмы, итераторы? Какова связь между ними?}
\begin{definition}[Контейнер]
Контейнер — это объект, который хранит набор элементов и предоставляет к ним доступ.
\end{definition}

\begin{definition}[Итератор]
Итератор — это абстракция, указывающая на элемент контейнера и позволяющая переходить между элементами.
\end{definition}

\begin{definition}[Алгоритм]
Алгоритм стандартной библиотеки — это обобщённая функция, которая работает с диапазонами элементов через итераторы.
\end{definition}

Связь между ними:
\begin{itemize}
    \item контейнер хранит данные;
    \item итератор даёт доступ к данным;
    \item алгоритм обрабатывает данные через итераторы.
\end{itemize}

\begin{remark}
Именно итераторы делают возможным принцип: один алгоритм — множество реализаций контейнеров.
\end{remark}

\subsection{Какая парадигма программирования описывает концепцию: один алгоритм — множество реализаций?}
Это \textbf{generic programming} — обобщённое, или шаблонное, программирование.

\begin{remark}
В C++ эта парадигма реализована через шаблоны: один и тот же алгоритм может работать с разными типами данных и контейнерами, если они удовлетворяют нужному интерфейсу.
\end{remark}

\subsection{В каком стиле стандартные алгоритмы применяются к стандартным контейнерам: ООП или функциональном? Почему?}
Стандартные алгоритмы применяются в \textbf{функциональном} стиле.

\begin{itemize}
    \item алгоритм — это отдельная функция;
    \item контейнер передаётся как диапазон через итераторы;
    \item результат зависит от входных данных, а не от скрытого состояния объекта.
\end{itemize}

\begin{example}
\begin{verbatim}
std::sort(v.begin(), v.end());
\end{verbatim}
\end{example}

Здесь алгоритм \texttt{sort} не является методом \texttt{vector}. Он работает с диапазоном, заданным итераторами.

\subsection{Чем схожи и чем различаются линейные и ассоциативные контейнеры?}
\begin{definition}[Линейные контейнеры]
Линейные контейнеры хранят элементы в последовательности.
\end{definition}

\begin{definition}[Ассоциативные контейнеры]
Ассоциативные контейнеры организуют доступ к элементам по ключу или по правилу упорядочения/хеширования.
\end{definition}

\textbf{Сходства}:
\begin{itemize}
    \item оба вида контейнеров хранят множество объектов;
    \item оба поддерживают стандартные операции доступа и обхода.
\end{itemize}

\textbf{Различия}:
\begin{itemize}
    \item линейные контейнеры ориентированы на порядок следования элементов;
    \item ассоциативные — на поиск по ключу или по свойству упорядоченности;
    \item у линейных контейнеров обычно проще доступ по позиции;
    \item у ассоциативных — лучше поиск по смысловому признаку.
\end{itemize}

\subsection{Чем схожи и чем различаются ассоциативные упорядоченные и неупорядоченные контейнеры?}
\textbf{Упорядоченные} ассоциативные контейнеры поддерживают элементы в некотором порядке, обычно по сравнению ключей.  
\textbf{Неупорядоченные} используют хеш-таблицы.

\textbf{Сходства}:
\begin{itemize}
    \item оба вида обеспечивают доступ по ключу;
    \item оба предназначены для быстрого поиска.
\end{itemize}

\textbf{Различия}:
\begin{itemize}
    \item упорядоченные контейнеры сохраняют порядок по ключу;
    \item неупорядоченные работают через hash function;
    \item упорядоченные обычно обеспечивают логарифмический поиск;
    \item неупорядоченные — в среднем близки к константному.
\end{itemize}

\subsection{\texttt{std::vector}: для чего применяется и в чём его назначение?}
\begin{definition}[std::vector]
\texttt{std::vector} — это динамически расширяемый массив, который хранит элементы подряд в памяти.
\end{definition}

Назначение:
\begin{itemize}
    \item хранить последовательность однотипных объектов;
    \item обеспечивать быстрый доступ по индексу;
    \item удобно наращиваться по мере добавления элементов.
\end{itemize}

\begin{remark}
\texttt{vector} — базовый универсальный линейный контейнер STL.
\end{remark}

\subsection{\texttt{std::vector}: доступ к элементам через \texttt{operator[]} и \texttt{at()}}
\texttt{operator[]} даёт доступ по индексу без проверки границ.  
\texttt{at()} делает доступ с проверкой границ и при ошибке выбрасывает исключение.

\begin{itemize}
    \item \texttt{v[i]} — быстрее, но потенциально опасно;
    \item \texttt{v.at(i)} — безопаснее, но чуть медленнее.
\end{itemize}

\begin{example}
\begin{verbatim}
std::vector<int> v{1, 2, 3};

int a = v[1];    // без проверки
int b = v.at(1); // с проверкой
\end{verbatim}
\end{example}

\subsection{\texttt{std::vector}: что возвращает \texttt{operator[]}?}
\texttt{operator[]} возвращает \textbf{ссылку} на элемент:
\begin{itemize}
    \item для неконстантного вектора — \texttt{T\&};
    \item для константного — \texttt{const T\&}.
\end{itemize}

\begin{remark}
Механизм именно ссылочный, потому что это доступ к уже существующему элементу внутри контейнера, а не копия элемента.
\end{remark}

\subsection{\texttt{std::vector}: варианты инициализации}
Основные варианты:
\begin{verbatim}
std::vector<int> a;
std::vector<int> b{1, 2, 3};
std::vector<int> c(5);        // 5 элементов по умолчанию
std::vector<int> d(5, 42);    // 5 элементов со значением 42
\end{verbatim}

Также можно инициализировать из другого диапазона:
\begin{verbatim}
std::vector<int> e(v.begin(), v.end());
\end{verbatim}

\subsection{\texttt{std::vector}: добавление и удаление элементов. \texttt{push\_back()}, \texttt{pop\_back()}, \texttt{insert()}, \texttt{erase()}}
\begin{itemize}
    \item \texttt{push\_back()} — добавление в конец;
    \item \texttt{pop\_back()} — удаление последнего элемента;
    \item \texttt{insert()} — вставка в произвольную позицию;
    \item \texttt{erase()} — удаление из произвольной позиции.
\end{itemize}

По эффективности:
\begin{itemize}
    \item \texttt{push\_back()} в среднем амортизированно быстро;
    \item \texttt{pop\_back()} — быстро;
    \item \texttt{insert()} и \texttt{erase()} в середине обычно медленнее, потому что требуется сдвиг элементов.
\end{itemize}

\begin{remark}
Самые дешёвые операции для \texttt{vector} — доступ по индексу и добавление/удаление с конца.
\end{remark}

\subsection{\texttt{std::vector}: распределение памяти, \texttt{reserve()}, \texttt{size()}, \texttt{capacity()}}
\begin{itemize}
    \item \texttt{size()} — число реально хранимых элементов;
    \item \texttt{capacity()} — объём выделенной памяти, доступный без новой аллокации;
    \item \texttt{reserve(n)} — заранее выделить память под минимум \texttt{n} элементов.
\end{itemize}

\begin{example}
\begin{verbatim}
std::vector<int> v;
v.reserve(1000);
\end{verbatim}
\end{example}

Это полезно, если заранее известно примерное число элементов.

\begin{remark}
Зарезервировать память при конструировании можно не всегда напрямую через отдельный параметр конструктора, но можно сразу создать нужный размер или затем вызвать \texttt{reserve()}.
\end{remark}

\subsection{\texttt{std::vector}: тип элементов и шаблонный параметр}
Тип элементов задаётся шаблонным параметром:
\begin{verbatim}
std::vector<int>
std::vector<std::string>
std::vector<MyType>
\end{verbatim}

После создания объекта изменить тип элементов уже нельзя: тип вектора фиксируется на этапе компиляции.

\begin{remark}
\texttt{std::vector<int>} и \texttt{std::vector<double>} — это разные типы.
\end{remark}

\subsection{\texttt{std::vector}: ограничения на тип элементов}
Вектор может хранить типы, которые:
\begin{itemize}
    \item можно размещать в памяти;
    \item можно копировать или перемещать при росте контейнера;
    \item можно разрушать;
    \item для некоторых конструкторов и операций — default constructable.
\end{itemize}

\begin{definition}[Default constructable]
Тип называется \textit{default constructable}, если его можно создать без аргументов.
\end{definition}

\begin{definition}[Copyable]
Тип называется \textit{copyable}, если его можно копировать.
\end{definition}

\begin{remark}
Современный \texttt{vector} использует и перемещение, но в учебной модели удобно помнить: элемент должен быть совместим с копированием/перемещением и разрушением.
\end{remark}

\subsection{Можно ли создать \texttt{std::vector<std::istream>}?}
Нет, обычный \texttt{std::vector<std::istream>} создать нельзя.

\begin{itemize}
    \item \texttt{std::istream} — неперемещаемый и некопируемый тип потока;
    \item потоками обычно не владеют как значениями в контейнере;
    \item вместо этого используют указатели, ссылки, \texttt{std::reference\_wrapper} или контейнеры умных указателей.
\end{itemize}

\begin{example}
\begin{verbatim}
std::vector<std::istream*> streams;
\end{verbatim}
\end{example}

\subsection{Можно ли создать вектор векторов? Чем ограничена вложенность?}
Да, можно:
\begin{verbatim}
std::vector<std::vector<int>> m;
std::vector<std::vector<std::vector<int>>> t;
\end{verbatim}

Вложенность ограничена:
\begin{itemize}
    \item памятью;
    \item удобством работы;
    \item требованиями к типам элементов на каждом уровне;
    \item глубиной, которую разумно поддерживать в программе.
\end{itemize}

\subsection{Как \texttt{vector} хранит элементы: по значению, по ссылке или через указатель?}
\texttt{std::vector} хранит элементы \textbf{по значению} внутри своего непрерывного буфера.

\begin{remark}
Он не хранит «ссылки на исходные элементы». При вставке элементы могут быть перемещены при перераспределении памяти.
\end{remark}

\subsection{\texttt{std::vector} и \texttt{std::array}: сходства и различия}
\textbf{Сходства}:
\begin{itemize}
    \item оба контейнера хранят элементы подряд;
    \item оба обеспечивают быстрый доступ по индексу;
    \item оба поддерживают работу с итераторами.
\end{itemize}

\textbf{Различия}:
\begin{itemize}
    \item \texttt{std::array} имеет фиксированный размер, известный на этапе компиляции;
    \item \texttt{std::vector} динамически растёт;
    \item \texttt{std::array} — обёртка над статическим массивом;
    \item \texttt{vector} управляет динамической памятью.
\end{itemize}

\subsection{\texttt{std::array} и C-массивы: сходства и различия}
\textbf{Сходства}:
\begin{itemize}
    \item фиксированный размер;
    \item хранение элементов подряд в памяти;
    \item быстрый доступ по индексу.
\end{itemize}

\textbf{Различия}:
\begin{itemize}
    \item \texttt{std::array} — полноценный объект со стандартным интерфейсом;
    \item C-массив при многих операциях деградирует до указателя;
    \item у \texttt{std::array} есть методы \texttt{size()}, \texttt{begin()}, \texttt{end()};
    \item C-массив не имеет встроенного безопасного интерфейса.
\end{itemize}

\begin{remark}
\texttt{std::array} обычно предпочтительнее C-массива в современном C++: он безопаснее и лучше интегрируется со стандартной библиотекой.
\end{remark}
